% GENERATED FILE. Edit canonical lessons, then run npm run generate:books.
% canonical-source-hash: fnv1a64:ac132940030589f6
% canonical-lessons: IT-C265-secondo-me, IT-C265-cosi, IT-C265-ripasso-opinione

\chapter{Secondo Me --- An Opinion, and What Follows}
\label{ch:it-an-opinion-and-what-follows}
\hlchaptermodality{265}

\begin{chapteropening}
\textbf{By the end of this chapter:} \emph{I can give my opinion with secondo me or credo di sì, and run a reason back with perché or forward with quindi and così.}

Complete IT-C265-ripasso-opinione: an opinion, its reason and its consequence, said in either order.
\end{chapteropening}

\section[secondo me]{secondo me — following me, which is to say, in my opinion}

\label{lesson:IT-C265-secondo-me}



\begin{quote}
Say \textquotedblleft{}the second\textquotedblright{}. (\emph{Il secondo.}) Then \textquotedblleft{}I think that\textquotedblright{}. (\emph{Penso che.})

A story says what happened. Giving reasons starts one step earlier, with what you think about it.
\end{quote}

\subsection*{You'll want to know: secondo me}
\textbf{secondo me} — \textquotedblleft{}in my opinion\textquotedblright{}, \textquotedblleft{}if you ask me\textquotedblright{}.

\begin{quote}
\textbf{Secondo me} il treno è in ritardo. — If you ask me, the train is late. \textbf{Secondo me} è vero. — In my opinion it's true.
\end{quote}

The \emph{me} is the stressed pronoun of \emph{per me}, and the others follow it: \textbf{secondo te?}, \textquotedblleft{}what do you think?\textquotedblright{}, is the easiest way to hand the opinion back to the other person.

And after \emph{secondo me} the verb stays plain. Chapter thirty-one warned that careful Italian bends the verb after \emph{penso che} into the \emph{congiuntivo}. \emph{Secondo me} asks for nothing of the kind: what follows is an ordinary sentence. That makes it the safest opinion-opener you own.

\begin{grammarlens}[title={Grammar Lens: credo di sì}]
You have \emph{credere}, \textquotedblleft{}to believe\textquotedblright{}. Its \textquotedblleft{}I\textquotedblright{} form is \textbf{credo}, and two short answers built on it are worth more than the verb:

\begin{quote}
\textbf{Credo di sì.} — I think so. \textbf{Credo di no.} — I don't think so.
\end{quote}

The \emph{di} has no English twin; Italian will not drop it. \emph{Credo che …} also exists and behaves exactly like \emph{penso che}, careful \emph{congiuntivo} and all.
\end{grammarlens}

\begin{cousinweb}
\textbf{secondo} is inherited from Latin \emph{secundum}, \textquotedblleft{}following, according to\textquotedblright{} (EtymDB 2.1), from the verb \emph{sequī}, \textquotedblleft{}to follow\textquotedblright{} (Lewis \& Short, under \emph{secundus}). The ordinal \textbf{secondo} is the same word: the second is the one that follows the first. \emph{Secondo me} is \textquotedblleft{}following me\textquotedblright{}, the view as it looks from where I stand. English \textbf{second}, \textbf{sequence} and \textbf{consequence} share it.

\textbf{credo} comes from Latin \emph{crēdere}, \textquotedblleft{}to trust, to believe\textquotedblright{} (EtymDB 2.1).
\end{cousinweb}

\subsection*{Your turn}
Say these aloud:

\begin{itemize}
\raggedright
  \item \emph{Secondo me il treno è in ritardo.}
  \item hand it back — \emph{Secondo te?}
  \item answer \emph{Il treno è in ritardo?} both ways — \emph{Credo di sì.} then \emph{Credo di no.}
  \item \emph{secondo}, the ordinal, then \emph{secondo me} — one word, two jobs
\end{itemize}

\begin{tcolorbox}[breakable,colback=teal!4,colframe=teal!35!black,title={Before you move on}]
What does \emph{secondo me} say, word for word? (\textbf{Following me.}) Why is it safer than \emph{penso che}? (The verb after it \textbf{stays plain}.) Say \textquotedblleft{}I think so\textquotedblright{}. (\emph{\textbf{Credo di sì.}}) And which Latin verb is behind both \emph{secondo} and \emph{sequence}? (\emph{\textbf{Sequī}}, \textquotedblleft{}to follow\textquotedblright{}.)
\end{tcolorbox}
\section[quindi, così]{quindi, così — the arrow that points forward}

\label{lesson:IT-C265-cosi}



\begin{quote}
Give an opinion with a reason: \textquotedblleft{}If you ask me, it's going to rain, because there's no sun.\textquotedblright{} (\emph{Secondo me piove, perché non c'è il sole.})

\emph{Perché} runs from the claim back to its reason. This lesson runs the other way.
\end{quote}

\begin{grammarlens}[title={Grammar Lens: the two arrows}]
The same two facts, in opposite orders:

\begin{quote}
Ho preso l'autobus \textbf{perché} pioveva. — I took the bus \textbf{because} it was raining. Pioveva, \textbf{quindi} ho preso l'autobus. — It was raining, \textbf{so} I took the bus.
\end{quote}

\textbf{Perché} points \textbf{back}, from what happened to its cause. \textbf{Quindi} points \textbf{forward}, from the cause to what followed. Which one you reach for depends only on which fact you want to say first, and that is a choice about emphasis, not about grammar.

Notice the tenses keep their jobs. \emph{Pioveva} is the scene, the reason; \emph{ho preso} is the event, the consequence.
\end{grammarlens}

\begin{grammarlens}[title={Grammar Lens: così, the spoken so}]
\textbf{così} can stand in the same place: \emph{Pioveva, \textbf{così} ho preso l'autobus.} It is lighter and more conversational than \emph{quindi}, which suits a letter or a report as well as speech. Both are right; \emph{così} is what you will hear most at a table, \emph{quindi} what you will read most.

You have been saying this word since chapter two. \textbf{Così così}, \textquotedblleft{}so-so\textquotedblright{}, is \emph{così} twice: \textquotedblleft{}this way, that way\textquotedblright{}.
\end{grammarlens}

\begin{cousinweb}
\textbf{così} comes from Latin \emph{eccum sīc}, \textquotedblleft{}look, thus\textquotedblright{} (EtymDB 2.1 lists it as a compound of the two). The \emph{sīc}, \textquotedblleft{}thus, in this way\textquotedblright{}, is the one English borrowed whole as \emph{sic}, the mark an editor writes after a quoted mistake: \textquotedblleft{}thus it stood\textquotedblright{}. \textbf{Quindi} has no row in EtymDB 2.1. The usual account takes it back to \emph{eccum inde}, \textquotedblleft{}look, from there\textquotedblright{}, and Latin \emph{inde} did mean \textquotedblleft{}from that place\textquotedblright{} (Lewis \& Short): from there, and so, therefore.
\end{cousinweb}

\subsection*{Your turn}
Say these aloud:

\begin{itemize}
\raggedright
  \item back — \emph{Ho preso l'autobus perché pioveva.}
  \item forward — \emph{Pioveva, quindi ho preso l'autobus.}
  \item the same, at a table — \emph{Pioveva, così ho preso l'autobus.}
  \item \emph{Secondo me piove, quindi resto a casa.}
\end{itemize}

\begin{tcolorbox}[breakable,colback=teal!4,colframe=teal!35!black,title={Before you move on}]
Which way does \emph{perché} point, and which way \emph{quindi}? (\emph{Perché} \textbf{back} to the cause; \emph{quindi} \textbf{forward} to the consequence.) Which of \emph{quindi} and \emph{così} is the lighter, spoken one? (\emph{\textbf{Così.}}) What did Latin \emph{sīc} mean? (\textbf{Thus}, in this way.)
\end{tcolorbox}
\section[(review)]{Review — what you think, why, and what follows}

\label{lesson:IT-C265-ripasso-opinione}



\begin{quote}
Someone asks \emph{Il museo è interessante?} Answer \textquotedblleft{}I think so\textquotedblright{}. (\emph{Credo di sì.}) Then ask them back. (\emph{Secondo te?})
\end{quote}

\begin{grammarlens}[title={What you've built}]
Three moves, and you own a word for each. An \textbf{opinion}: \emph{secondo me}. A \textbf{reason}, pointing back: \emph{perché}. A \textbf{consequence}, pointing forward: \emph{quindi} or \emph{così}.

\begin{quote}
\emph{Secondo me il museo è interessante, perché è molto antico.} \emph{Il treno è caro, quindi prendo l'autobus.}
\end{quote}

The first sentence runs claim, then cause. The second runs cause, then consequence. Put them side by side and you can argue a small point in either direction, which is most of what giving reasons is.
\end{grammarlens}

\subsection*{Your turn}
Say these aloud:

\begin{itemize}
\raggedright
  \item \emph{Secondo me il museo è interessante, perché è molto antico.}
  \item the same facts, forward — \emph{Il museo è molto antico, quindi è interessante.}
  \item \emph{Il treno è caro, così prendo l'autobus.}
  \item a reason in the past — \emph{Pioveva, quindi sono restato a casa.}
  \item your own opinion about a place you know, with a \emph{perché}
\end{itemize}

\begin{tcolorbox}[breakable,colback=teal!4,colframe=teal!35!black,title={Before you move on}]
Name the three moves and a word for each. (Opinion \emph{\textbf{secondo me}}, reason \emph{\textbf{perché}}, consequence \emph{\textbf{quindi}} or \emph{\textbf{così}}.) Say \textquotedblleft{}I don't think so\textquotedblright{}. (\emph{\textbf{Credo di no.}}) In \emph{pioveva, quindi sono restato a casa}, which half is the scene? (\emph{\textbf{Pioveva}} — the reason.)
\end{tcolorbox}
